\chapter{Desarrollo, implantación y validación}
\label{sec:desarrollo_implantacion_validacion}

\section{Organización del desarrollo}
\subsection{Metodología de trabajo}
\subsection{Gestión del proyecto en GitHub Projects}
\subsection{Plan de pruebas}

La validación del proyecto se organizó de forma incremental, siguiendo la misma división por sprints utilizada durante el desarrollo. El objetivo fue comprobar, en cada etapa, que las funcionalidades incorporadas no solo estaban disponibles desde la interfaz, sino que respetaban las reglas de negocio, los permisos, los contratos de datos y las restricciones definidas para el sistema.

La estrategia de pruebas combinó dos enfoques complementarios. Por un lado, se implementaron pruebas automatizadas sobre las partes más críticas del proyecto: esquemas compartidos, servicios de la API y determinadas utilidades del portal web. Por otro lado, los recorridos completos de usuario, especialmente aquellos relacionados con la aplicación móvil, el uso offline, la cámara, la ubicación y la sincronización, se validaron mediante pruebas manuales funcionales apoyadas en evidencias visuales.

Esta separación permite reflejar con precisión el alcance real de la validación. La memoria no atribuye al proyecto pruebas automatizadas que no existen, pero sí deja constancia de que las funcionalidades principales fueron comprobadas mediante el tipo de prueba más adecuado a su naturaleza.

\begin{table}[H]
\centering
\caption{Estrategia general de pruebas del proyecto}
\label{tab:estrategia_general_pruebas}
\small
\begin{tabularx}{\textwidth}{p{3.2cm}X X}
\toprule
Tipo de prueba & Objetivo & Aplicación en OmniLitter \\
\midrule
Pruebas de contrato & Verificar que los datos intercambiados entre clientes y API cumplen la estructura esperada. & Esquemas compartidos de autenticación, campañas, catálogo, sesiones, observaciones y fotografías. \\
Pruebas unitarias de servicio & Comprobar reglas de negocio sin depender de la interfaz ni de una base de datos real. & Servicios de usuarios, grupos, campañas, catálogo y observaciones del backend mediante mocks de Prisma. \\
Pruebas automatizadas web & Validar lógica auxiliar del portal web que no depende directamente de componentes visuales. & Cálculo de distancia de rutas y preferencias de exportación del dashboard. \\
Pruebas manuales funcionales & Comprobar recorridos completos desde la perspectiva del usuario final. & Portal web, aplicación móvil, captura de campo, uso offline, sincronización y evidencias visuales. \\
Pruebas de regresión & Comprobar que los cambios posteriores no rompen funcionalidades ya validadas. & Reejecución de la suite automatizada y revisión manual de los flujos principales al cierre de cada sprint. \\
\bottomrule
\end{tabularx}
\end{table}

La suite automatizada se ejecutó con \textit{Vitest}. En el paquete compartido se utilizaron esquemas definidos con \textit{Zod}, mientras que en el backend se validaron servicios mediante mocks de Prisma. Esta decisión permitió comprobar reglas de negocio relevantes sin depender de una base de datos PostgreSQL real ni de la ejecución completa de rutas HTTP.

\begin{table}[H]
\centering
\caption{Alcance real de las pruebas automatizadas}
\label{tab:alcance_real_pruebas_automatizadas}
\small
\begin{tabularx}{\textwidth}{p{2.4cm}p{3.4cm}p{2.4cm}X}
\toprule
Área & Ficheros implementados & Herramienta & Cobertura principal \\
\midrule
API & Tests de \texttt{apps/api} & Vitest & Servicios de campañas, catálogo, grupos, usuarios y observaciones con mocks de Prisma. \\
Web & Tests de \texttt{apps/web} & Vitest & Preferencias de exportación del dashboard y cálculo de distancia de rutas. \\
Shared & Test de \texttt{packages/shared} & Vitest + Zod & Esquemas compartidos de autenticación, campañas, catálogo, sesiones, observaciones y fotografías. \\
Mobile & Sin ficheros de test configurados & No aplica & La aplicación móvil se validó mediante pruebas manuales funcionales y evidencias visuales. \\
\bottomrule
\end{tabularx}
\end{table}

Para facilitar la trazabilidad, las pruebas automatizadas se agruparon mediante códigos identificativos. Estos códigos se emplean en los apartados de cada sprint para relacionar las tareas desarrolladas con las comprobaciones realizadas.

\begin{table}[H]
\centering
\caption{Leyenda de códigos usados en las pruebas}
\label{tab:leyenda_codigos_pruebas}
\small
\begin{tabularx}{\textwidth}{p{2.7cm}X}
\toprule
Código & Significado \\
\midrule
SCH & Pruebas de contrato sobre esquemas compartidos. Validan que los datos enviados por web, móvil o API cumplen la estructura esperada. \\
USR & Pruebas del servicio de usuarios del backend. \\
GRP & Pruebas del servicio de grupos y membresías del backend. \\
CMP & Pruebas del servicio de campañas del backend. \\
CAT & Pruebas del catálogo científico usado por el sistema. \\
OBS & Pruebas del servicio de observaciones del backend. \\
WEB-SES & Pruebas de utilidades web relacionadas con sesiones y cálculo de distancias. \\
WEB-EXP & Pruebas de preferencias de exportación en el portal web. \\
\bottomrule
\end{tabularx}
\end{table}

La ejecución global de la suite automatizada finalizó correctamente. En total se ejecutaron 77 pruebas automatizadas distribuidas entre el paquete compartido, el portal web y la API.

\begin{table}[H]
\centering
\caption{Resultado global de la ejecución de pruebas automatizadas}
\label{tab:resultado_global_pruebas_automatizadas}
\small
\begin{tabularx}{\textwidth}{p{4cm}p{4cm}r r}
\toprule
Paquete & Resultado & Ficheros & Casos \\
\midrule
\texttt{@omnilitter/shared} & Correcto & 1 & 23 \\
\texttt{@omnilitter/web} & Correcto & 2 & 15 \\
\texttt{@omnilitter/api} & Correcto & 5 & 39 \\
\texttt{@omnilitter/mobile} & Sin tests automatizados configurados & 0 & 0 \\
\midrule
Total ejecutado & Correcto & 8 & 77 \\
\bottomrule
\end{tabularx}
\end{table}

\subsection{Criterios de finalización de tareas}
\subsection{Estructura de releases}

\clearpage
\section{Sprint 1: Base funcional del sistema}
\subsection{Objetivo del sprint}
\subsection{Parte del proyecto abordada}
\subsection{Tiempo invertido y reparto del trabajo}
\subsection{Implementación realizada}
\subsection{Pruebas realizadas}

En el primer sprint se validó la base funcional del sistema. Las pruebas se centraron en autenticación, administración inicial de usuarios y grupos, campañas, catálogo, sesiones de muestreo y primeras observaciones. Esta fase era especialmente importante porque establecía los contratos y servicios sobre los que se apoyarían los sprints posteriores.

\subsubsection{Pruebas automatizadas implementadas}

\begin{table}[H]
\centering
\caption{Pruebas automatizadas vinculadas al Sprint 1}
\label{tab:pruebas_automatizadas_sprint1}
\scriptsize
\begin{tabularx}{\textwidth}{p{3.5cm}p{3cm}X p{1.8cm}}
\toprule
Tarea o funcionalidad & Cobertura automatizada & Comprobación principal & Resultado \\
\midrule
Login, registro y cierre de sesión & SCH-01, SCH-02 & Validación del registro con contraseña robusta, confirmación coincidente y rechazo de contraseñas débiles o confirmaciones incorrectas. & Correcto \\
Portal de administración de usuarios & USR-01, USR-03 & Bloqueo de correos duplicados y creación de usuarios con o sin membresías iniciales. & Correcto \\
Gestión de grupos y membresías & GRP-06, GRP-07 & Restauración de grupos, permisos de administración y rechazo de acciones realizadas por usuarios sin autorización. & Correcto \\
Campañas y sesiones backend & SCH-04 a SCH-13, CMP-01 & Estados válidos de campañas, referencias de catálogo y reglas de geometría para sesiones libres, por cuadrícula o por transecto. & Correcto \\
Catálogo científico & CAT-01 a CAT-08 & Creación, actualización, detección de duplicados, borrado y restricciones por uso de entornos y tipologías. & Correcto \\
API de observaciones y adjuntos & SCH-14 a SCH-23, OBS-04 a OBS-06 & Observaciones con ítems obligatorios, coordenadas válidas, cantidades positivas, tipos de imagen soportados y control de acceso a observaciones propias o ajenas. & Correcto \\
Formulario móvil de observación & SCH-14 a SCH-16 & Validación compartida de sesión, taxonomía, coordenadas, ítems y cantidades utilizada por el formulario móvil. & Correcto \\
\bottomrule
\end{tabularx}
\end{table}

Estas pruebas protegen reglas que no deben depender únicamente de la interfaz. Aunque los formularios ayuden a evitar datos incorrectos, los esquemas compartidos y los servicios de backend garantizan que la validación se mantiene ante datos reutilizados por distintos clientes o enviados directamente a la API.

\subsubsection{Pruebas manuales funcionales}

Las pruebas manuales de este sprint se reservan para los recorridos completos de usuario, como registro, inicio de sesión, administración básica, creación de campañas, creación de sesiones y registro inicial de observaciones desde la aplicación móvil. Su detalle se documenta mediante evidencias visuales en el apartado correspondiente de validación manual.

\subsection{Desviaciones e incidencias}
\subsection{Resultado del sprint}

\clearpage
\section{Sprint 2: Captura de campo y análisis inicial}
\subsection{Objetivo del sprint}
\subsection{Parte del proyecto abordada}
\subsection{Tiempo invertido y reparto del trabajo}
\subsection{Implementación realizada}
\subsection{Pruebas realizadas}

El segundo sprint reforzó la captura de datos de campo y el análisis inicial. La validación se orientó a comprobar la calidad del dato capturado, la incorporación de fotografías, el uso de ubicación GPS, el control de permisos sobre observaciones y las primeras funcionalidades de análisis del portal web.

\subsubsection{Pruebas automatizadas implementadas}

\begin{table}[H]
\centering
\caption{Pruebas automatizadas vinculadas al Sprint 2}
\label{tab:pruebas_automatizadas_sprint2}
\scriptsize
\begin{tabularx}{\textwidth}{p{3.5cm}p{3cm}X p{1.8cm}}
\toprule
Tarea o funcionalidad & Cobertura automatizada & Comprobación principal & Resultado \\
\midrule
GPS en sesión y observación & WEB-SES-01 a WEB-SES-03 & Cálculo de distancia para rutas vacías, rutas con un único punto y suma de distancia Haversine entre puntos consecutivos. & Correcto \\
Fotografías por observación y metadatos & SCH-17 a SCH-23 & Aceptación de formatos de imagen soportados y rechazo de tipos MIME no permitidos. & Correcto \\
Validación científica de observaciones & SCH-14 a SCH-16 & Existencia de ítems, coordenadas válidas y cantidades positivas antes de aceptar una observación. & Correcto \\
Permisos de observaciones & OBS-01 a OBS-06 & Alcance de observaciones por usuario, visibilidad ampliada para gestores y bloqueo de observaciones ajenas sin permisos. & Correcto \\
Dashboard de métricas científicas & WEB-EXP-01 a WEB-EXP-12 & Formatos soportados, persistencia de preferencias en \texttt{localStorage} y recuperación ante valores obsoletos. & Correcto \\
\bottomrule
\end{tabularx}
\end{table}

La cobertura de este sprint resulta relevante para la calidad científica del sistema. Una observación no debe almacenarse sin ítems, con coordenadas imposibles, con cantidades no positivas o con ficheros de imagen no soportados. Además, las pruebas de permisos aseguran que cada usuario consulta únicamente la información que corresponde a su ámbito de autorización.

\subsubsection{Pruebas manuales funcionales}

Las pruebas manuales de este sprint cubren los flujos que requieren interacción real con la aplicación móvil y revisión visual del resultado: permisos de ubicación, captura de coordenadas, asociación de fotografías, validación de formularios, consulta por roles y visualización inicial de métricas en el portal web.

\subsection{Desviaciones e incidencias}
\subsection{Resultado del sprint}

\clearpage
\section{Sprint 3: Sincronización, mapas y exportación científica}
\subsection{Objetivo del sprint}
\subsection{Parte del proyecto abordada}
\subsection{Tiempo invertido y reparto del trabajo}
\subsection{Implementación realizada}
\subsection{Pruebas realizadas}

El tercer sprint estuvo orientado a sincronización, mapas y exportación científica. En esta fase fue necesario diferenciar con claridad entre las pruebas automatizadas existentes y la validación funcional de integración. La sincronización móvil-backend se comprobó mediante recorridos manuales completos, pero no se implementó una prueba automática de integración contra una base de datos real.

\subsubsection{Pruebas automatizadas implementadas}

\begin{table}[H]
\centering
\caption{Pruebas automatizadas vinculadas al Sprint 3}
\label{tab:pruebas_automatizadas_sprint3}
\scriptsize
\begin{tabularx}{\textwidth}{p{3.5cm}p{3cm}X p{1.8cm}}
\toprule
Tarea o funcionalidad & Cobertura automatizada & Comprobación principal & Resultado \\
\midrule
Exportación CSV y soporte de rutas & WEB-SES-01 a WEB-SES-03 & Cálculo de distancia de tracks como soporte para visualización, análisis de recorridos y exportación de información asociada a sesiones. & Correcto \\
Preferencias de exportación científica & WEB-EXP-01 a WEB-EXP-12 & Formatos admitidos, persistencia de la preferencia de exportación y recuperación ante formatos no soportados. & Correcto \\
Configuración y cambio de contraseña & SCH-03, USR-04 a USR-06 & Rechazo de reutilización de contraseña, usuario inexistente, contraseña actual incorrecta y guardado seguro del nuevo hash. & Correcto \\
Sincronización móvil-backend & Sin suite automatizada específica & La sincronización se validó mediante pruebas manuales de flujo completo. No se atribuye cobertura automatizada inexistente. & Manual \\
\bottomrule
\end{tabularx}
\end{table}

La ausencia de una prueba automatizada específica para sincronización se considera una limitación de la validación automática. No obstante, se documenta de forma explícita para mantener la trazabilidad y evitar que la memoria presente como automatizada una comprobación que realmente fue funcional y manual.

\subsubsection{Pruebas manuales funcionales}

Las pruebas manuales de este sprint son especialmente relevantes, ya que varias funcionalidades dependen de la interacción entre aplicación móvil, API y portal web. Incluyen trabajo sin conexión, sincronización posterior, visualización en mapas, exportación científica y configuración de usuario.

\subsection{Desviaciones e incidencias}
\subsection{Resultado del sprint}

\clearpage
\section{Sprint 4: Cierre funcional y refinamiento}
\subsection{Objetivo del sprint}
\subsection{Parte del proyecto abordada}
\subsection{Tiempo invertido y reparto del trabajo}
\subsection{Implementación realizada}
\subsection{Pruebas realizadas}

El cuarto sprint correspondió al cierre funcional y al refinamiento del sistema. La validación tuvo un carácter principalmente regresivo: se reejecutaron las pruebas automatizadas existentes y se revisaron manualmente los flujos principales del portal web y de la aplicación móvil para comprobar que los cambios finales no introducían errores.

\subsubsection{Pruebas automatizadas implementadas}

\begin{table}[H]
\centering
\caption{Pruebas automatizadas vinculadas al Sprint 4}
\label{tab:pruebas_automatizadas_sprint4}
\scriptsize
\begin{tabularx}{\textwidth}{p{3.5cm}p{3cm}X p{1.8cm}}
\toprule
Tarea o funcionalidad & Cobertura automatizada & Comprobación principal & Resultado \\
\midrule
Página de configuración web & USR-02, USR-04 a USR-06, GRP-01 a GRP-15, CMP-02 a CMP-04, OBS-01 a OBS-06 & Permisos por rol, vista agregada de superadministrador, participación en campañas, archivado y restauración de grupos, membresías y alcance de observaciones. & Correcto \\
Entorno en sesión y rangos científicos & CAT-01 a CAT-08, SCH-08 a SCH-13 & Validación de entornos, tipologías, referencias de catálogo y reglas compartidas que condicionan campañas y sesiones. & Correcto \\
Cierre funcional y refinamiento & Suite completa & Reejecución de las pruebas de API, web y shared para comprobar que los ajustes finales no rompen funcionalidades previamente validadas. & Correcto \\
Aplicación móvil & Sin suite automatizada específica & Los ajustes visuales y funcionales de la aplicación móvil se validaron mediante recorridos manuales. & Manual \\
\bottomrule
\end{tabularx}
\end{table}

Esta validación permitió comprobar que las reglas principales del sistema seguían funcionando tras los cambios finales. Se mantiene como limitación que la aplicación móvil no dispone de una suite automatizada configurada.

\subsubsection{Pruebas manuales funcionales}

Las pruebas manuales de este sprint se orientan a la revisión global del sistema: autenticación, configuración, campañas, sesiones, observaciones, catálogo, navegación del portal web y recorrido principal de la aplicación móvil.

\subsection{Desviaciones e incidencias}
\subsection{Resultado del sprint}

\clearpage
\section{Resumen global del desarrollo}
\subsection{Funcionalidades implementadas}
\subsection{Evolución del sistema}
\subsection{Cobertura de requisitos}
\subsection{Matriz de trazabilidad prueba/requisito}
\subsection{Limitaciones finales}
